\section{The elliptic deformation and exclusion of the floor}
\label{n3:deformation:section}

We define the coupled system, derive its scalar-curvature identity, and
choose a continuation path along which the conformal factor cannot touch
its lower bound.  We use the fixed data of
Proposition~\ref{n3:domains:prepared-domain}. Thus $\Omega$ has one controlled
asymptotically flat end, $K$ is compactly supported, and its compact boundary
$B=B_b\cup B_w$ consists of black and white faces. Either type may be absent.
The normal $\nu$ points into $\Omega$, and
\[
 H=\Div_B\nu,\qquad P_B=\tr_BK.
\]
The numbers $b_-<0<b_+$, the smooth compactly supported offset $C$, and
the positive function $\rho=O(r^{-4})$ are fixed as in that proposition.
In particular,
\begin{equation}\label{n3:deformation:face-thresholds}
 b_-+C=P_B+H\ \hbox{on }B_b,\qquad
 b_++C=P_B-H\ \hbox{on }B_w.
\end{equation}
All compact sets and geometric constants in this section are fixed after
this preparation. Extend the end radius to a smooth function on
$\overline\Omega$ bounded below by one. Write $\Omega_R$ for the large
spherical truncations and $S_R$ for their outer boundary.

A \emph{polynomial constant} is bounded by $C(1+n)^N$, with $C,N$
depending on the prepared data and fixed auxiliary profiles, but not on
$R$, the homotopy parameter, or the solution. Constants permitted to
depend arbitrarily on fixed $n$ will be identified explicitly.
First fix $0<\epsilon<1$, then choose $n$ sufficiently large, and finally
require $R\ge R_{\min}(n,\epsilon)$. We always take $n\ge4$ and
$n>2/\epsilon$.

\subsection{Variables and equations}
\label{n3:deformation:variables}

Choose a smooth nonincreasing $\vartheta:\R\to[0,1]$ equal to one on
$(-\infty,0]$ and zero on $[1,\infty)$. Define
\begin{equation}\label{n3:deformation:l-profile}
 \begin{gathered}
 p(t)=n\vartheta(nt),\qquad
 l(t)=\exp\left(\int_0^t p(s)\dd s\right),\\
 L(t)=l(t)e^{2t},\qquad p_L=p+2.
 \end{gathered}
\end{equation}
Thus $l(0)=1$, $l(t)=e^{nt}$ for $t\le0$, and $l$ is constant for
$t\ge1/n$. Derivatives of $p$ of every fixed order are polynomially bounded.
Set
\begin{equation}\label{n3:deformation:small-parameters}
 \ell=e^{-\epsilon n},\qquad \tau=\ell^{3/2}.
\end{equation}
Whenever $t\ge-\epsilon$, one has $\ell\le l(t)\le e$.

For a function $f$, set
\begin{equation}\label{n3:deformation:basic-variables}
 \begin{gathered}
 \sigma=|\nabla f|,\quad D=1+l^2\sigma^2,\quad d=D^{-1},\\
 a=l\sigma/\sqrt D,\quad v=a^2=1-d,\quad
 w=l\nabla f/\sqrt D,\quad \chi=4d/(4+pv),\\
 Z=t+\tfrac18\log D,\quad u=L\sqrt D,\quad U=l\sqrt D,\quad h=\tau f.
 \end{gathered}
\end{equation}
Derivatives and contractions refer to $g$ unless indicated otherwise.
At $\sigma>0$, write $e=\nabla f/\sigma$ and
$A_j=I+(j-1)e\otimes e$. The matrices used below have the direction-free
expressions
\begin{equation}\label{n3:deformation:axial-matrices}
 A_d=I-w\otimes w,\qquad
 A_\chi=I-\frac{p+4}{4+pv}w\otimes w.
\end{equation}
Vectors and covectors are identified with $g$.

The unknowns will be $(f,Z)$. For fixed $\nabla f$,
$\partial Z/\partial t=1+pv/4>0$, and its defining expression tends to
$-\infty$ and $+\infty$ at the two ends of the $t$-axis. It therefore
defines a unique smooth $t=t(Z,\nabla f)$ for all inputs, without a floor
assumption. The formulas above are smooth at $\nabla f=0$, and
$0<\chi\le d\le1$ for finite inputs.
Define
\begin{equation}\label{n3:deformation:metrics}
 \bar g=g+l^2\dd f^2,\qquad \hat g=e^{4t}\bar g,\qquad
 H^f=\frac l{\sqrt D}\Hess f,\qquad S=K+H^f.
\end{equation}
The original trace equation is
\begin{equation}\label{n3:deformation:trace}
 \tr_{A_\chi}S=F,\qquad F=h+C,
\end{equation}
where $\tr_A S$ is contraction with the contravariant matrix $A$.
Along the continuation path we will specify modified right sides $F$.
The full flux is
\begin{equation}\label{n3:deformation:flux}
 V=uA_\chi\bigl(4\nabla Z+K(w,\cdot)\bigr).
\end{equation}
At $S_R$ prescribe $f=Z=0$. Initially the inner Dirichlet values are
$h=b_-$ on $B_b$ and $h=b_+$ on $B_w$.

\subsection{The exact scalar-curvature identity}

The next identity separates the original constraint densities, a
nonnegative quadratic term, and a divergence.  The second equation will
prescribe that divergence so that the deformed metric has nonnegative
scalar curvature.

\begin{proposition}[Deformation identity]\label{n3:deformation:identity}
Let $f,t$ be smooth and $F=\tr_{A_\chi}(K+H^f)$. At $\sigma>0$ decompose
$S$ relative to $\R e\oplus e^\perp$:
\[
 T=S|_{e^\perp\times e^\perp},\quad
 M=S(e,\cdot)|_{e^\perp},\quad q=S(e,e),\quad \tr T=F-\chi q.
\]
Then
\begin{equation}\label{n3:deformation:scalar-identity}
 \begin{split}
 \tfrac12e^{4t}\Scal_{\hat g}
 ={}&8\pi(\mu+J(w))+\mathcal T
       +w(F)-F\tr K-u^{-1}\Div V,
 \end{split}
\end{equation}
where
\begin{equation}\label{n3:deformation:quadratic-term}
 \begin{split}
 \mathcal T={}&\tfrac12|T^{\operatorname{tf}}|^2
 +(M+pa\nabla_\perp t)\cdot(M+(p+2)a\nabla_\perp t)\\
 &+4(p+1)|\dd t|_{A_d}^2
   +(\chi-\chi^2/4)q^2+2(p+1)\chi qa\,\partial_e t\\
 &-\chi Fq/2+3F^2/4.
 \end{split}
\end{equation}
After substituting $F=\tr_{A_\chi}S$, this expression extends smoothly
across $\sigma=0$.
\end{proposition}

\begin{proof}
Introduce an auxiliary coordinate $s_0$ and the stationary Lorentz metric
\[
 {\bf g}=-l^2(\dd s_0-\dd f)^2+\bar g
       =-U^2\dd s_0^2+g(\dd x+\beta\dd s_0,\dd x+\beta\dd s_0),
 \qquad \beta=l^2\nabla f=Uw.
\]
Its $s_0$-slices have future normal $N_0=(\partial_{s_0}-\beta)/U$ and
second fundamental form
\begin{equation}\label{n3:deformation:auxiliary-k}
 k=-\frac{\sym\nabla\beta}{U}
   =-H^f-p(\dd t\otimes w^\flat+w^\flat\otimes\dd t).
\end{equation}
Write $\theta=\tr k$ and, for symmetric tensors, put
\[
 \mathscr B(A,B)=\langle A,B\rangle-(\tr A)(\tr B),\qquad
 P_A=A-(\tr A)g,\qquad Q=K-k.
\]
The normal expansion equation and contracted Gauss equation are
\[
 \begin{split}
 \Ric_{\bf g}(N_0,N_0)&=-N_0(\theta)-|k|^2+U^{-1}\Delta U,\\
 \Scal_g&=\Scal_{\bf g}+2\Ric_{\bf g}(N_0,N_0)+|k|^2-\theta^2.
 \end{split}
\]
Consequently
$\Scal_{\bf g}=\Scal_g+|k|^2+\theta^2+
2N_0(\theta)-2U^{-1}\Delta U$.
Stationarity gives $UN_0(\theta)=-\beta(\theta)$, and
$\Div\beta=-U\theta$ follows from \eqref{n3:deformation:auxiliary-k}.
Thus
\begin{equation}\label{n3:deformation:stationary-scalar}
 U\Scal_{\bf g}
 =U\bigl(\Scal_g+\mathscr B(k,k)\bigr)
   -2\Div(\nabla U+\theta\beta).
\end{equation}
In the coordinates $s_0-f$ the same metric is a static warped product.
Its mixed Christoffel symbols are
$\mathbf{\Gamma}^{s}_{is}=\partial_i\log l$ and
$\mathbf{\Gamma}^{i}_{ss}=l\bar\nabla^il$; contraction gives
$\Scal_{\bf g}=\Scal_{\bar g}-2l^{-1}\bar\Delta l$.
Since $\bar g^{-1}=A_d$ and $\dd V_{\bar g}=(U/l)\dd V_g$,
\begin{equation}\label{n3:deformation:bar-laplacian}
 \bar\Delta\phi=\frac lU\Div\left(\frac Ul A_d\nabla\phi\right).
\end{equation}
Direct differentiation of $U$ and \eqref{n3:deformation:auxiliary-k} gives
\[
 \nabla U-\frac Ul A_d\nabla l=-k(\beta,\cdot).
\]
In fact both sides equal
\[
 \frac{l^4}{U}\Hess f(\nabla f,\cdot)
 +\frac{l^3}{U}\bigl[|\nabla f|^2\dd l
                       +(\nabla f)(l)\dd f\bigr],
\]
as follows by differentiating $U^2=l^2+l^4|\nabla f|^2$.
Substitution in \eqref{n3:deformation:stationary-scalar} yields
\[
 U\Scal_{\bar g}
 =U\bigl(\Scal_g+\mathscr B(k,k)\bigr)+2\Div(P_k\beta).
\]
The constraints read
$16\pi\mu=\Scal_g-\mathscr B(K,K)$ and $8\pi J=\Div P_K$.
Using $P_K:\nabla\beta=-U\mathscr B(K,k)$, expansion therefore gives
\begin{equation}\label{n3:deformation:unweighted-identity}
 U\Scal_{\bar g}
 =U[16\pi(\mu+J(w))+\mathscr B(Q,Q)]-2\Div(P_Q\beta).
\end{equation}
Indeed the derivative of $P_K$ cancels its term in $\Div(P_K\beta)$,
and
$\mathscr B(Q,Q)+2\mathscr B(K,k)=\mathscr B(K,K)+\mathscr B(k,k)$.

The three-dimensional conformal formula is
\[
 \tfrac12e^{4t}\Scal_{\hat g}
 =\tfrac12\Scal_{\bar g}-4\bar\Delta t-4|\dd t|_{A_d}^2.
\]
For $Y=P_Qw$, replacing $U$ by $u=e^{2t}U$ gives
$-U^{-1}\Div(UY)=-u^{-1}\Div(uY)+2Y(t)$.
Since $u/(U/l)=L$, Equation~\eqref{n3:deformation:bar-laplacian} similarly gives
\[
 -4\bar\Delta t-4|\dd t|_{A_d}^2
 =-u^{-1}\Div(4uA_d\nabla t)+4(p+1)|\dd t|_{A_d}^2.
\]
Before the flux modification, the scalar-curvature expression is
\begin{equation}\label{n3:deformation:pre-flux}
 \begin{split}
 8\pi(\mu+J(w))+\tfrac12\mathscr B(Q,Q)
 +2P_Q(w,\nabla t)+4(p+1)|\dd t|_{A_d}^2\\
 -u^{-1}\Div\bigl(u(P_Qw+4A_d\nabla t)\bigr).
 \end{split}
\end{equation}

At $\sigma>0$,
\begin{equation}\label{n3:deformation:Z-differential}
 4\dd Z=(4+pv)\dd t+v\dd\log\sigma.
\end{equation}
Write $x=\partial_e t$ and $y=\nabla_\perp t$. The blocks of $Q$ are
\[
 \begin{gathered}
 Q_{\perp\perp}=T,\qquad Q_{e\perp}=M+pay,\qquad Q_{ee}=q+2pax,\\
 \tr Q=F+(1-\chi)q+2pax.
 \end{gathered}
\]
Hence $(P_Qw)_\perp=a(M+pay)$ and $(P_Qw)_e=a(\chi q-F)$.
Together with $aH^f(e,\cdot)=v\dd\log\sigma$,
Equation~\eqref{n3:deformation:Z-differential} proves component by component
\begin{equation}\label{n3:deformation:flux-rewrite}
 V/u=P_Qw+4A_d\nabla t+wF.
\end{equation}
Also $uw=e^{2t}l^2\nabla f$, so
\begin{equation}\label{n3:deformation:uw-divergence}
 \Div(uw)/u
 =F+(1-\chi)q-\tr K+2(p+1)ax.
\end{equation}
Replacing the flux in \eqref{n3:deformation:pre-flux} adds
$w(F)+F\Div(uw)/u$. The remaining quadratic expression is
\begin{equation}\label{n3:deformation:quadratic-invariant}
 \begin{split}
 \mathcal T={}&\tfrac12\mathscr B(Q,Q)+2P_Q(w,\nabla t)
                  +4(p+1)|\dd t|_{A_d}^2\\
 &+F\left[F+\frac{p+4}{4+pv}S(w,w)+2(p+1)w(t)\right].
 \end{split}
\end{equation}
Here $(1-\chi)q=(p+4)S(w,w)/(4+pv)$,
which gives a smooth formula at $w=0$.

For the explicit coefficients, substitute
$|T|^2=|T^{\operatorname{tf}}|^2+(F-\chi q)^2/2$ in
$\mathscr B(Q,Q)/2$.
The transverse terms become $(M+pay)\cdot(M+(p+2)ay)$.
The coefficients of $q^2,Fq,F^2$ are respectively
$\chi-\chi^2/4,-\chi/2,3/4$.
The $qx$ coefficient is $2(p+1)\chi a$, the $Fx$ terms cancel, and
the remaining $x^2$ term is $4(p+1)dx^2$.
These are the terms in \eqref{n3:deformation:quadratic-term}.
\end{proof}

\begin{lemma}[Polynomial coercivity]\label{n3:deformation:coercivity}
The expression $\mathcal T$ is nonnegative. For $0\le p\le n$,
\begin{equation}\label{n3:deformation:coercive-estimate}
 |\dd t|_{A_d}^2+|T|^2+|M|^2+dq^2+F^2
 \le C(1+n)^N\mathcal T
\end{equation}
with universal $C,N$.
At $\dd t=0$ there is the stronger comparison, uniform in $p,d$,
\begin{equation}\label{n3:deformation:floor-norm}
 c\bigl(|T^{\operatorname{tf}}|^2+|M|^2+\chi q^2+F^2\bigr)
 \le\mathcal T
 \le C\bigl(|T^{\operatorname{tf}}|^2+|M|^2+\chi q^2+F^2\bigr).
\end{equation}
\end{lemma}

\begin{proof}
With $x=\partial_e t$ and $y=\nabla_\perp t$, exact completion gives
\begin{equation}\label{n3:deformation:square-completion}
 \begin{split}
 \mathcal T={}&\tfrac12|T^{\operatorname{tf}}|^2
 +|M+(p+1)ay|^2+[4(p+1)-v]|y|^2\\
 &+4(p+1)d\left(x+\frac{\chi aq}{4d}\right)^2\\
 &+\tfrac34\left(F-\frac{\chi q}{3}\right)^2
 +\frac{4d(9-d)}{3(4+pv)^2}q^2.
 \end{split}
\end{equation}
Completing first in $x$ leaves $3\chi q^2/(4+pv)$; completing next in
$F$ leaves $3\chi/(4+pv)-\chi^2/12=4d(9-d)/[3(4+pv)^2]$.
The latter coefficient is at least $32d/[3(4+n)^2]$.
It controls $dq^2$ polynomially. The $F$-square then controls $F^2$,
since $\chi^2q^2\le dq^2$.
The transverse gradient coefficient is at least three, and the shifted
$M$-square controls $M$ with an additional factor $(1+n)^2$.
The squared $d$-norm of the shift in the $x$-square is
$\chi^2vq^2/(16d)\le dq^2/16$, so it also controls $dx^2$.
Finally $|T|^2=|T^{\operatorname{tf}}|^2+(F-\chi q)^2/2$.
This proves \eqref{n3:deformation:coercive-estimate}.
If $\dd t=0$, the sole cross term is $-\chi Fq/2$.
Using $|\chi Fq|/2\le\chi q^2/4+\chi F^2/4$ and $0<\chi\le1$
proves \eqref{n3:deformation:floor-norm}.
\end{proof}

\begin{lemma}[Boundary identity]\label{n3:deformation:boundary-identity}
On a face where $f$ is constant,
\begin{equation}\label{n3:deformation:boundary}
 V_\nu/u=d(H+4\partial_\nu t)-H+w_\nu(F-P_B).
\end{equation}
Its mean curvature in $\hat g$, with normal into $\Omega$, is
$e^{-2t}\sqrt d\,(H+4\partial_\nu t)$.
\end{lemma}

\begin{proof}
The tangential Hessian is
$\Hess f|_{TB}=(\partial_\nu f)\operatorname{II}_g$.
Thus $\tr_\perp H^f=w_\nu H$ and
$\chi q=F-P_B-w_\nu H$.
The normal component of \eqref{n3:deformation:flux}, using
\eqref{n3:deformation:Z-differential}, is
\[
 V_\nu/u=4d\partial_\nu t+\chi w_\nu q
        =4d\partial_\nu t+w_\nu(F-P_B)-vH.
\]
The $\bar g$-normal is $\sqrt d\,\nu$. Its mean curvature is $\sqrt d\,H$,
because the tangential metric is unchanged and its normal first variation
is multiplied by $\sqrt d$. The conformal mean-curvature formula gives
the asserted mean curvature for $\hat g$.
\end{proof}

\subsection{The source, boundary condition, and continuation path}

We now prescribe the divergence and inner flux. For the original system,
the source is chosen to give the deformed metric nonnegative scalar
curvature. A penalty near the lower bound will exclude contact with that
bound along the continuation path.

Choose $3/4<\gamma<1$ and fixed positive smooth weights
\[
 \rho_0\asymp r^{-4},\qquad \rho_1\asymp r^{-2\gamma},
\]
with differentiated end bounds. Let $m:\R\to[0,1]$ be smooth, equal to one
on $(-\infty,0]$ and zero on $[1,\infty)$, and set
\[
 m_0(t)=m(n(t+\epsilon)),\qquad
 \mathcal P=C_n(\rho_0+v\rho_1).
\]
The positive $\delta_0$ is fixed independently of $n,R$ in
Proposition~\ref{n3:deformation:floor}; thereafter $C_n$ is a sufficiently
large polynomial constant. Define
\begin{equation}\label{n3:deformation:source}
 \Xi=\delta_0\mathcal T+\rho+\delta_0\tau a\sigma-m_0(t)\mathcal P.
\end{equation}
Choose a smooth face function $N_B>1+\sup_B|H|$.
The original second equation and its inner condition are
\begin{equation}\label{n3:deformation:original-second-equation}
 \Div V=u\Xi,\qquad V_\nu/u=-H+w_\nu(F-P_B)-N_Bd.
\end{equation}
Lemma~\ref{n3:deformation:boundary-identity} gives
$H+4\partial_\nu t=-N_B$.
Only $-\epsilon\le t\le-\epsilon+1/n$ is penalized in the admissible range.
Since $n>2/\epsilon$, this band has negative $t$, and
\begin{equation}\label{n3:deformation:penalty-band}
 \ell\le l\le e\ell,\qquad c\ell\le L\le C\ell
 \quad\hbox{on the admissible penalty band}.
\end{equation}
The latter constants can be uniform for $0<\epsilon<1$.

\begin{proposition}[Specified homotopy]\label{n3:deformation:homotopy}
There is a continuous piecewise smooth family of systems indexed by
$s\in[0,4]$, beginning at
\eqref{n3:deformation:trace}, \eqref{n3:deformation:original-second-equation}
and the original Dirichlet values, and ending at $f=0$ and the scalar
problem with zero source and zero inner conormal flux.
Throughout its first three stages the source is
\eqref{n3:deformation:source}, with $\mathcal T$ evaluated using the actual
trace $F$. Each trace equation has right-side derivative at least one
in $h$, with $x,Z,\nabla f$ fixed.
\end{proposition}

\begin{proof}
Fix a smooth $j:\R\to[0,1]$ which vanishes for $t\le0$ and equals one
for $t\ge1$. For $0\le\lambda\le1$ put
\[
 V_\lambda=uA_\chi(4\nabla Z+\lambda K(w,\cdot)),\qquad
 \mathcal B=-H+w_\nu(F-P_B)-N_Bd.
\]
Before the last stage use the boundary condition
\begin{equation}\label{n3:deformation:homotopy-boundary}
 \begin{split}
 (V_\lambda)_\nu/u
 ={}&[1-(1-\lambda)j(t)]\\
 &\cdot[\mathcal B-(1-\lambda)(A_\chi K(w,\cdot))_\nu].
 \end{split}
\end{equation}
Outer values remain $f=Z=0$ throughout.

We specify two smooth collar modifications. Let
$\nu_s=\nabla s/|\nabla s|$ be the outward unit normal to a collar leaf;
the leaf coordinate $s$ here is local and is separate from the homotopy
parameter. Choose disjoint collar cutoffs equal to one at the faces.
Take $\eta>0$ so small that $b_-+3\eta<b_+$.
A black contribution to $D_0(x,w,h)$ is a negative constant times a
collar cutoff, times a smooth directional cutoff supported where
$w\cdot\nu_s<-1/2$ and equal to one at $-1$, times a nonincreasing
height cutoff equal to one below $b_-+\eta$ and zero above $b_-+2\eta$.
A white contribution has positive sign, its directional cutoff supported
where $w\cdot\nu_s>1/2$ and equal to one at $1$, and a nondecreasing
height cutoff equal to zero below $b_+-2\eta$ and one above $b_+-\eta$.
Decrease $\eta$ to separate these ranges. Their sum $D_0$ is smooth,
bounded, supported in the collars, and nondecreasing in $h$.
It is nonpositive for $h\le b_-+\eta$, nonnegative for
$h\ge b_+-\eta$, and vanishes on black outward and white inward
directions, for every length $0\le|w|\le1$.

Choose the magnitudes larger than $2\sup_B|H|+1$.
Equation~\eqref{n3:deformation:face-thresholds} then gives, for
$F=h+C+D_0$ at the assigned boundary height $b=b_\pm$,
\begin{equation}\label{n3:deformation:directional-barriers}
 F(x,+\nu,b)\ge P_B+H,\qquad F(x,-\nu,b)\le P_B-H.
\end{equation}
These are limiting directional evaluations at $a=1,\chi=0$.
The compatible direction has equality; $D_0$ corrects the other one.
Its derivatives on bounded height ranges are independent of $n$.
Choose a smooth compactly supported extension $b(x)$, constant equal
to $b_\pm$ on each smaller collar.
On these collars let $D_1(x,w)=A_1w\cdot\nu_s$, extended by a collar
cutoff, with $A_1>\sup_B|H|+1$.
Then $D_1(x,0)=0$ and $D_1(x,\pm\nu)=\pm A_1$ on $B$.

The four stages are the following.
\begin{enumerate}[label=\textup{(\roman*)}]
 \item For $0\le s\le1$, set $\lambda=1-s$, keep $F=h+C$ and $h|_B=b$,
 and use $\Div V_\lambda=u\Xi$ and
 \eqref{n3:deformation:homotopy-boundary}.
 \item For $1\le s\le2$, set $\lambda=0$, $\alpha=s-1$, retain $h|_B=b$,
 and replace the trace right side by $F=h+C+\alpha D_0(x,w,h)$.
 Keep the same scalar equation and boundary prescription.
 \item For $2\le s\le3$, set $\lambda=0$, $\zeta=s-2$, prescribe
 $h|_B=(1-\zeta)b$, and use
 \begin{equation}\label{n3:deformation:interpolation-trace}
 \begin{split}
 F={}&h+(1-\zeta)[C+D_0(x,w,h+\zeta b(x))]\\
    &+\zeta[\tr_{A_\chi}K+D_1(x,w)].
 \end{split}
 \end{equation}
 Keep the same scalar equation and boundary prescription.
 At a face, $h+\zeta b=b$. The value at $w=\pm\nu$ is the convex
 combination of the preceding directional value and
 $P_B+D_1(x,\pm\nu)$, so \eqref{n3:deformation:directional-barriers} persists.
 \item For $3\le s\le4$, retain the trace problem at $\zeta=1$ and
 multiply both $\Xi$ and the right side of
 \eqref{n3:deformation:homotopy-boundary} by $\eta_s=4-s$.
\end{enumerate}
The prescriptions agree at the common endpoints.
Their $h$-derivatives are $1$, $1+\alpha\partial_hD_0$, or
$1+(1-\zeta)\partial_hD_0$, all at least one.

At $\zeta=1$ the trace equation is
$\tr_{A_\chi}H^f=h+D_1(x,w)$ with zero Dirichlet values.
At a positive interior maximum of $f$, one has $w=0$, $A_\chi=I$,
$D_1=0$, and $l\Delta f=\tau f>0$, a contradiction.
A negative minimum is excluded in the same way.
Thus $f=0$ and $t=Z$ in the last stage.
At $s=4$ the remaining equation is
$\Div(4L(Z)\nabla Z)=0$, with zero inner conormal flux and $Z=0$ on
$S_R$. Multiplication by $Z$ and integration give $Z=0$.
For the fixed-point construction in Section~\ref{n3:existence:section},
the last map therefore has output identically zero.
\end{proof}

\subsection{Elementary height and end barriers}

\begin{lemma}[Height control]\label{n3:deformation:height}
Every smooth trace solution in Proposition~\ref{n3:deformation:homotopy}
satisfies $|h|\le C_0$, independently of $n,R$, the homotopy parameter,
and the input $Z$. Its actual right side satisfies $|F|\le C_F$
with the same independence.
\end{lemma}

\begin{proof}
The collar modifications vanish at $w=0$.
At an interior extremum the trace equation reads either
$\tr K+l\Delta f=h+C$ or
$\tr K+l\Delta f=h+(1-\zeta)C+\zeta\tr K$.
The maximum and minimum inequalities bound $h$ by fixed bounds for
$|\tr K|+|C|$. Boundary values range between $b_-,0,b_+$.
This proves the height bound without using the floor.
Boundedness of $D_0,D_1$ and $|\tr_{A_\chi}K|\le3|K|$ proves the
assertion for $F$.
\end{proof}

\begin{lemma}[End height and outer boundary]\label{n3:deformation:end-height}
There are $r_0,C>0$, independent of $n,R$ and the homotopy parameter, such
that for $R\ge2r_0$,
\begin{equation}\label{n3:deformation:end-height-estimate}
 |h|\le C(r^{-\gamma}-R^{-\gamma})\quad(r_0\le r\le R).
\end{equation}
In particular,
\begin{equation}\label{n3:deformation:outer-gradient}
 |\nabla f|\le C R^{-1-\gamma}/\tau\quad\hbox{on }S_R.
\end{equation}
For fixed $\epsilon,n$, a sufficiently large lower bound for $R$ gives
$t>-\epsilon$ on $S_R$.
\end{lemma}

\begin{proof}
Take $r_0$ outside the supports of $K,C,D_0,D_1$.
There $F=h$ in the first three stages, and $h=0$ in the last one.
For $\psi=r^{-\gamma}$, with axis $\nabla\psi/|\nabla\psi|$,
the controlled end gives, uniformly for $0<\chi\le1$,
\[
 \tr_{A_\chi}\Hess\psi
 =\gamma((\gamma+1)\chi-2)r^{-\gamma-2}
      +O(r^{-\gamma-3})<0
\]
after enlarging $r_0$.
At a comparison contact, coefficients use the given $Z$ and the common
gradient of test and solution. The positive factor $l/(\tau\sqrt D)$
preserves this sign. Height monotonicity therefore makes
$C(\psi-R^{-\gamma})$ an upper barrier and its negative a lower barrier.
Choose $C$ to dominate $C_0$ on $r=r_0$ for every $R\ge2r_0$.
Differentiating at the common outer value gives
\eqref{n3:deformation:outer-gradient}; tangential derivatives there vanish.

At $S_R$, $Z=0$. The monotonicity of its defining expression implies
$t>-\epsilon$ if $\ell^2|\nabla f|^2<e^{8\epsilon}-1$.
It suffices to require
$C^2\ell^2\tau^{-2}R^{-2-2\gamma}<e^{8\epsilon}-1$,
which is a permissible lower bound depending on fixed $n,\epsilon$.
\end{proof}

\subsection{Exclusion of a contact with the floor}

We exclude a proposed floor contact using the equations at that point.
These estimates require neither an upper bound for $Z$ nor a bound for
$|\nabla f|$.

\begin{lemma}[Curvature at a minimum of $t$]\label{n3:deformation:floor-curvature}
At an interior point where $\dd t=0$ and $\Hess t\ge0$,
\begin{equation}\label{n3:deformation:floor-gauss}
 \tfrac12e^{4t}\Scal_{\hat g}
 \le \tfrac12\Scal_g-v\Ric_g(e,e)+\mathcal S(H^f),
\end{equation}
where, for a symmetric tensor $A$,
\begin{equation}\label{n3:deformation:gauss-quadratic}
 \begin{split}
 \mathcal S(A)={}&\tfrac14(\tr_\perp A)^2
 -\tfrac12|A_{\perp\perp}^{\operatorname{tf}}|^2\\
 &+dA_{ee}\tr_\perp A-d|A_{e\perp}|^2.
 \end{split}
\end{equation}
There are universal $c_*,\eta_*>0$ and polynomial constants $C_n$ such
that, at $\dd t=0$,
\begin{equation}\label{n3:deformation:floor-algebra-general}
 \mathcal T-\mathcal S(H^f)
 \ge c_*\mathcal T-C_n(F^2+|K|^2).
\end{equation}
If $K=0$, the more useful uniform estimate is
\begin{equation}\label{n3:deformation:floor-algebra-end}
 \mathcal T-\mathcal S(H^f)\ge c_*\mathcal T-C_n vF^2.
\end{equation}
\end{lemma}

\begin{proof}
Consider the graph of $f$ in the Riemannian product with warped fiber
$(\Omega\times\R,g+l^2\dd b_0^2)$. Its induced metric is $\bar g$.
At $\dd t=0$ one also has $\dd l=0$, and its second fundamental form,
up to orientation, is $H^f$. Write $E=l^{-1}\partial_{b_0}$.
The graph unit normal is $\sqrt d\,E-ae$. The ambient curvature formulas
at this point are
\[
 \begin{gathered}
 \Scal_{\mathrm{amb}}=\Scal_g-2\Delta l/l,\\
 \Ric_{\mathrm{amb}}|_{T\Omega}=\Ric_g-\Hess l/l,\qquad
 \Ric_{\mathrm{amb}}(E,E)=-\Delta l/l,
 \end{gathered}
\]
with zero mixed Ricci terms. Hence the ambient contribution to half the
graph scalar curvature is
\[
 \tfrac12\Scal_g-v\Ric_g(e,e)
      -v\,\tr_\perp\Hess l/l.
\]
The last term is nonpositive: at $\dd t=0$,
$\Hess l=lp\,\Hess t$ and $p\ge0$.
The second-form contribution is
$[(\tr_{\bar g}H^f)^2-|H^f|_{\bar g}^2]/2$.
Since the inverse graph metric is $A_d$, its expansion is exactly
\eqref{n3:deformation:gauss-quadratic}.
Finally the conformal contribution at this point is
$-4\bar\Delta t\le0$. This proves \eqref{n3:deformation:floor-gauss}.

For the algebraic estimates we must retain the degenerate normal
coefficient. First evaluate $\mathcal S$ on
$S=K+H^f$, whose trace in the transverse plane is $F-\chi q$.
Subtracting \eqref{n3:deformation:gauss-quadratic} from
\eqref{n3:deformation:quadratic-term} at $\dd t=0$ gives
\begin{equation}\label{n3:deformation:floor-difference}
 \begin{split}
 \mathcal T-\mathcal S(S)
 ={}&|T^{\operatorname{tf}}|^2+(1+d)|M|^2\\
 &+\chi(1+d-\chi/2)q^2-dFq+F^2/2.
 \end{split}
\end{equation}
The normal coefficient is at least $\chi$, while
$|dFq|\le\chi q^2/2+d^2F^2/(2\chi)$.
Because
\[
 d^2/\chi=d(4+pv)/4\le1+n/4,
\]
Equation~\eqref{n3:deformation:floor-norm} implies
$\mathcal T-\mathcal S(S)\ge c\mathcal T-C(1+n)F^2$
with universal $c>0$.

To replace $S$ by $H^f=S-K$, polarize
\eqref{n3:deformation:gauss-quadratic}. The difference is bounded by
\[
 C|K|(|T^{\operatorname{tf}}|+|M|+|F|+d|q|)+C|K|^2.
\]
Indeed the only extra normal product is $dq\,\tr_\perp K$;
all remaining products involve transverse entries, $dM$, or
$\tr_\perp S=F-\chi q$.
By \eqref{n3:deformation:floor-norm} and $d/\sqrt\chi\le C\sqrt{1+n}$,
this bound is at most
$C\sqrt{1+n}|K|\sqrt{\mathcal T}+C|K|^2$.
Young's inequality absorbs a fixed small fraction of $\mathcal T$
and costs only $C(1+n)|K|^2$. This proves
\eqref{n3:deformation:floor-algebra-general}.

If $(1+p)v\le\eta_*$, then $d$ and $\chi$ are within $C\eta_*$
of one. At $d=\chi=1$ the normal block in
\eqref{n3:deformation:floor-difference} is
$3q^2/2-Fq+F^2/2$; its symmetric matrix has smallest eigenvalue
$(2-\sqrt2)/2>0$.
Choose $\eta_*>0$ small enough that its coefficients remain within
half this spectral gap. The transverse coefficients are already
bounded below. For $K=0$, Equation~\eqref{n3:deformation:floor-norm}
therefore gives $\mathcal T-\mathcal S(H^f)\ge c_*\mathcal T$
in this regime.
In the remaining regime $v>\eta_* /(1+p)$, use
\eqref{n3:deformation:floor-algebra-general} with $K=0$ and absorb
$(1+n)F^2$ into $C(1+n)^2vF^2$.
This proves \eqref{n3:deformation:floor-algebra-end}.
All formulas extend to $\sigma=0$ by their invariant expressions;
there $d=\chi=1$ and the full quadratic expressions are independent of
the temporary choice of axis.
\end{proof}

\begin{lemma}[Homotopy errors at a floor contact]
\label{n3:deformation:floor-errors}
There is a fixed compact set $\mathcal K\subset\overline\Omega$ containing
the supports of $K,C,b,D_0,D_1$ such that, at an interior minimum
$t=-\epsilon$ of a smooth solution in the first three stages,
\begin{equation}\label{n3:deformation:floor-divergence-lower}
 \Div V_\lambda/u
 \ge 2\delta_0\mathcal T+\tau a\sigma
      -C(1+n)^N
       [\mathbf1_{\mathcal K}+v(\rho_0+\rho_1)].
\end{equation}
Here $\delta_0>0$ can be fixed sufficiently small independently of
$n,R$ and the homotopy parameter.
\end{lemma}

\begin{proof}
At this contact $\dd t=0$ and $p=n$ is locally constant as a function
of $t$ near its value $-\epsilon$. Differentiating $w$ in an orthonormal
frame gives
\begin{equation}\label{n3:deformation:floor-w-derivative}
 \nabla_iw_j=H^f_{ij}-w_jH^f_{ik}w^k.
\end{equation}
In the $e$-frame the normal vector slot is multiplied by $d$.
Equations \eqref{n3:deformation:floor-norm} and $d/\sqrt\chi\le C\sqrt{1+n}$
therefore show
\begin{equation}\label{n3:deformation:floor-w-control}
 |\nabla w|\le C(1+n)^N(|K|+\sqrt{\mathcal T}).
\end{equation}
No bound for the uncontracted normal Hessian is asserted.

Each first-three-stage right side can be written as a smooth function
$F(x,w,h,p)$. Its $h$-derivative is at least one.
On the bounded height range supplied by Lemma~\ref{n3:deformation:height},
its explicit $x$- and $w$-derivatives are polynomially bounded and
compactly supported, apart from the derivative of the term $h$.
This follows directly from the fixed collar cutoffs and
$A_\chi=I-(p+4)w\otimes w/(4+pv)$, whose denominator is at least four.
For example its first $w$-derivatives on $|w|\le1$ are bounded by
$C(1+n)^2$, and the other explicit modifications have bounds independent
of $n$. Combining these bounds with
\eqref{n3:deformation:floor-w-control} and Young's inequality therefore
produces a fixed finite power of $1+n$.
Since $\nabla p=0$ at the contact and $w(h)=\tau a\sigma$, the chain rule
and \eqref{n3:deformation:floor-w-control} give, for every fixed $\eta>0$,
\begin{equation}\label{n3:deformation:floor-F-control}
 w(F)\ge\tau a\sigma-\eta\mathcal T
                 -C_\eta(1+n)^N\mathbf1_{\mathcal K}.
\end{equation}
The dependence $D_0(x,w,h+\zeta b(x))$ only adds the fixed compact
derivative of $b$ and does not change this conclusion.

The omitted flux is $(1-\lambda)uA_\chi K(w,\cdot)$.
Let $B_K=A_\chi K(w,\cdot)$. Its derivative is bounded by a polynomial
constant times $1+\sqrt{\mathcal T}$ on $\mathcal K$, by
\eqref{n3:deformation:floor-w-control}; derivatives of $p$ vanish.
The weight term requires the contraction rather than a bound for
$|\nabla\log u|$. At the contact,
\[
 \nabla\log u=H^f(w,\cdot),\qquad
 \langle\nabla\log u,B_K\rangle
 =a^2\bigl(\chi H^f_{ee}K_{ee}
                    +H^f_{e\perp}\cdot K_{e\perp}\bigr).
\]
This involves only $\chi H^f_{ee}$ and $H^f_{e\perp}$ and is controlled
by \eqref{n3:deformation:floor-norm}. It follows that
\begin{equation}\label{n3:deformation:floor-flux-control}
 |\Div(uB_K)|/u
 \le\eta\mathcal T+C_\eta(1+n)^N\mathbf1_{\mathcal K}.
\end{equation}

Combine the scalar identity with \eqref{n3:deformation:floor-gauss}.
The difference $8\pi\mu-\Scal_g/2$ equals
$((\tr K)^2-|K|^2)/2$ and is compactly supported; so is $J$.
The term $F\tr K$ is compactly bounded by Lemma~\ref{n3:deformation:height}.
On $\mathcal K$, Equation~\eqref{n3:deformation:floor-algebra-general}
therefore contributes $c_*\mathcal T-C(1+n)^N$.
Outside $\mathcal K$, $K=0$, $F=h$, and
\eqref{n3:deformation:floor-algebra-end} together with
Lemma~\ref{n3:deformation:end-height} gives
$c_*\mathcal T-C(1+n)^Nv r^{-2\gamma}$.
The remaining Ricci term is bounded below by
$-C\mathbf1_{\mathcal K}-Cv r^{-3}$; enlarge $\mathcal K$ if necessary.
Both end weights are bounded by a multiple of $\rho_0+\rho_1$.
Finally subtract the omitted flux and use
\eqref{n3:deformation:floor-F-control}--\eqref{n3:deformation:floor-flux-control},
choosing their two $\eta$'s with sum at most $c_*/2$.
Fix $0<\delta_0<\min(c_*/4,1/4)$. Decreasing it slightly if necessary
gives \eqref{n3:deformation:floor-divergence-lower}.
\end{proof}

\begin{proposition}[Floor exclusion]\label{n3:deformation:floor}
Choose $\delta_0$ as in Lemma~\ref{n3:deformation:floor-errors}.
There is a polynomial choice $C_n=C(1+n)^N$ in the penalty such that,
for all sufficiently large $n$ and all
$R\ge R_{\mathrm{floor}}(n,\epsilon)$, no smooth solution anywhere along
Proposition~\ref{n3:deformation:homotopy} can satisfy $t\ge-\epsilon$
and attain $t=-\epsilon$.
In particular every such admissible solution has $t>-\epsilon$ on
$\overline{\Omega_R}$.
\end{proposition}

\begin{proof}
Choose $R_{\mathrm{floor}}$ to include the requirement of
Lemma~\ref{n3:deformation:end-height}. A floor contact cannot occur on $S_R$.
In the first three stages, an inner contact has $j(t)=0$.
Equation~\eqref{n3:deformation:homotopy-boundary} then cancels the omitted
flux on both sides and gives $V_\nu/u=\mathcal B$.
Equation~\eqref{n3:deformation:boundary} forces
$H+4\partial_\nu t=-N_B$, so $\partial_\nu t<0$.
This contradicts the nonnegative derivative into $\Omega_R$ at a
boundary minimum.

Consider an interior contact. Let
$E_n=C(1+n)^N[\mathbf1_{\mathcal K}+v(\rho_0+\rho_1)]$
denote the error in \eqref{n3:deformation:floor-divergence-lower}.
Because $\rho_0$ has a positive minimum on $\mathcal K$,
$\mathbf1_{\mathcal K}\le C_{\mathcal K}\rho_0$.
Also $v\rho_0\le\rho_0$, and $\rho/\rho_0$ is globally bounded.
We can therefore choose a polynomial $C_n$ so large that, pointwise
for all $0\le v\le1$,
\begin{equation}\label{n3:deformation:penalty-choice}
 \mathcal P\ge E_n+2\rho+\rho_0.
\end{equation}
At the floor $m_0=1$, so the scalar equation and the lower estimate give
\[
 \begin{split}
 0
 &\ge \delta_0\mathcal T+(1-\delta_0)\tau a\sigma
                     +\mathcal P-E_n-\rho\\
 &\ge \rho+\rho_0>0,
 \end{split}
\]
a contradiction. Enlarging $C_n$ below preserves this inequality and
preserves its polynomial dependence.

In the last stage $f=0$, $t=Z$, $u=L$, and $v=0$.
At an interior minimum,
$\mathcal T=(|K|^2+(\tr K)^2)/2$.
Increase $C_n$ polynomially if necessary so that
$\delta_0\mathcal T+\rho-C_n\rho_0<0$ at all such points.
For $\eta_s>0$, the scalar equation at a floor minimum has left side
$4L\Delta t\ge0$ and strictly negative right side.
At an inner floor minimum its flux is
$4\partial_\nu t=\eta_s(-H-N_B)<0$, also impossible.
If $\eta_s=0$, multiplication of
$\Div(4L(t)\nabla t)=0$ by $t$, with zero outer value and zero inner
flux, gives $t=0$.
Thus the last stage has no floor contact either.
\end{proof}

We have excluded floor contact for every smooth admissible solution
along the specified homotopy. Section~\ref{n3:apriori:section} supplies the
remaining a priori bounds, and Section~\ref{n3:existence:section} proves
existence. The choice of $\mathcal P$ uses only polynomial losses in $n$;
it is independent of all fixed-$n$ Schauder constants.
